\question

El lema de la serpiente es una herramienta importante en álgebra
homológica, para la cual se utilizan diagramas conmutativos.
El diagrama que se muestra en la Figura se extiende
en el lema de la serpiente al llamado diagrama de serpiente.

\begin{figure}[ht!]
	\centering
	\begin{tikzpicture}[>=angle 90]
		\matrix[matrix of math nodes,
		row sep=3em, column sep=3em,,
		text height=1.5ex, text depth=0.25ex]
		{
		&&|[name=A1]| A_1 &|[name=B1]| B_1 &|[name=C1]| C_1 &|[name=01]| 0 \\
		&|[name=00]| 0 &|[name=A0]| A_0 &|[name=B0]| B_0 &|[name=C0]| C_0 \\
		};
		\draw[>->,font=\scriptsize,blue]
		(A0) edge node[auto] {$\varphi_1$} (B0);
		\draw[->>,font=\scriptsize,blue]
		(B1) edge node[auto] {$\psi_1$} (C1);
		\draw[->,font=\scriptsize,blue]
		(A1) edge node[auto] {$\varphi_1$} (B1)
		(B0) edge node[auto] {$\psi_0$} (C0)
		(A1) edge node[left] {$d$} (A0)
		(B1) edge node[auto] {$\delta$} (B0)
		(C1) edge node[auto] {$\partial$} (C0)
		(C1) edge (01)
		(00) edge (A0);
	\end{tikzpicture}
	\caption{Diagrama conmutativo con filas exactas.}
	\label{fig:1}
\end{figure}

\begin{figure}[ht!]
	\centering
	\begin{tikzpicture}[>=angle 90]
		\matrix[matrix of math nodes,
		row sep=3em, column sep=3em,,
		text height=1.5ex, text depth=0.25ex]
		{
		&&|[name=kd]| \Ker(d) & |[name=kdelta]| \Ker(\delta)
		&|[name=kdel]| \Ker(\partial) \\
		&&|[name=A1]| A_1 &|[name=B1]| B_1 &|[name=C1]| C_1 &|[name=01]| 0 \\
		&|[name=00]| 0 &|[name=A0]| A_0 &|[name=B0]| B_0 &|[name=C0]| C_0 \\
		&&|[name=cd]| \Coker(d) &|[name=cdelta]| \Coker(\delta) &|[name=cdel]|
		\Coker(\partial) \\
		};
		\draw[right hook->,font=\scriptsize]
		(kd) edge (A1)
		(kdelta) edge (B1)
		(kdel) edge (C1);
		\draw[>->,font=\scriptsize,blue]
		(A0) edge node[below] {$\varphi_0$} (B0);
		\draw[->>,font=\scriptsize,blue]
		(B1) edge node[auto] {$\psi_1$} (C1);
		\draw[->,font=\scriptsize,blue]
		(A1) edge node[auto] {$\varphi_1$} (B1)
		(B0) edge node[below] {$\psi_0$} (C0)
		(A1) edge node[left] {$d$} (A0)
		(B1) edge node[auto] {$\delta$} (B0)
		(C1) edge node[auto] {$\partial$} (C0)
		(C1) edge (01)
		(00) edge (A0);
		\draw[->>,font=\scriptsize]
		(A0) edge (cd)
		(B0) edge (cdelta)
		(C0) edge (cdel);
		\draw[->,red]
		(kd) edge (kdelta)
		(kdelta) edge (kdel)
		(cd) edge (cdelta)
		(cdelta) edge (cdel);
		\draw[->,dashed,red]
		(kdel) edge[out=0,in=180,red] node[above left] {$D$} (cd);
	\end{tikzpicture}
	\caption{Diagrama de la serpiente (El lema de la serpiente recibe
		su nombre de la flecha $D$, que serpentea a través del diagrama
		como una serpiente).}
	\label{fig:2}
\end{figure}

Dado que todo matemático debe saber cómo representar diagramas
conmutativos en texto, esto se practicará aquí utilizando el diagrama
de serpiente.

\begin{parts}
	\part

	Reproduzca el diagrama que se muestra en la Figura~\ref{fig:1} en
	\LaTeX.
	No es necesario que se ciña a las flechas, los colores ni la
	notación utilizada.
	Sin embargo, el contenido de los diagramas debe ser coherente.
	Asegúrese de que las expresiones matemáticas también estén
	configuradas en modo matemático dentro de la figura.
	Existen varios paquetes que pueden facilitarle esta tarea.
	Por ejemplo, los diagramas conmutativos se pueden crear de forma
	bastante elegante con TikZ.

	\part

	Establezca el lema de la serpiente completo, incluyendo la
	demostración.
	Utilice los paquetes \AmS{} y defina usted mismo los operadores y
	entornos necesarios mediante
	\mintinline{latex}|\DeclareMathOperator| y
	\mintinline{latex}|\newtheorem|.
	La representación de los elementos matemáticos en el texto y en las
	figuras (diagramas) debe ser, por supuesto, la misma.
\end{parts}

\begin{solution}
	\begin{tikzpicture}[xscale=0.5,yscale=0.5]
		\draw[help lines] (0,0) grid (32,16);
	\end{tikzpicture}
\end{solution}
